% Part: first-order-logic
% Chapter: introduction
% Section: syntax

\documentclass[../../../include/open-logic-section]{subfiles}

\begin{document}

\olfileid{fol}{int}{syn}

\olsection{विन्यासमीमांसा}

प्रथम-क्रम तर्कशास्त्रातील कोणत्या चिन्हमाला !!{sentence}s म्हणून गणल्या
जातात हे प्रथम काटेकोरपणे सांगावे लागेल. हे आपण नंतर करू; सध्या
उदाहरणांच्या साहाय्याने पुढे जाऊ. प्रथम-क्रम तर्कशास्त्राचे मूलभूत
रचनाघटक---म्हणजे त्याचा शब्दसंग्रह---दोन भागांत विभागलेले आहेत. पहिल्या
भागात विशिष्ट गोष्टींबद्दल काही सांगण्यासाठी किंवा विशिष्ट गोष्टी
निर्देशित करण्यासाठी आपण वापरतो ती चिन्हे येतात. गोष्टी निर्देशित
करण्यासाठी आपण !!{constant}s वापरतो आणि निर्देशित केलेल्या गोष्टींबद्दल
काही सांगण्यासाठी !!{predicate}s वापरतो. उदाहरणार्थ, एखादी एक गोष्ट
निर्देशित करण्यासाठी $\Obj a$ हे !!a{constant} म्हणून वापरू शकतो आणि मग
!!{sentence}~$\Atom{\Obj P}{\Obj a}$ वापरून तिच्याबद्दल काही सांगू शकतो.
``$\Obj a$'' आणि ``$\Obj P$'' यांचे काही अर्थ तुम्ही मनात धरले असतील,
तर $\Atom{\Obj P}{\Obj a}$ हे इंग्रजीतील वाक्य म्हणून वाचू शकता (आणि
आकारिक तर्कशास्त्र प्रथम शिकताना तुम्ही बहुधा तसे केलेही असेल).
प्रथम-क्रम तर्कशास्त्रातील अशी साधी !!{sentence}s मिळाल्यानंतर,
शब्दसंग्रहातील दुसरा भाग---तार्किक चिन्हे (संयोजक आणि संख्यापक)---वापरून
त्यांपासून अधिक जटिल वाक्ये रचता येतात. म्हणून, उदाहरणार्थ,
$(\Atom{\Obj P}{\Obj a} \land \Atom{\Obj Q}{\Obj b})$ किंवा
~$\lexists[\Obj x][\Atom{\Obj P}{\Obj x}]$ यांसारख्या अभिव्यक्ती
रचता येतात.

काटेकोर अधितार्किक अभ्यासासाठी आवश्यक असलेल्या चिन्हार्थमीमांसेच्या
नेमक्या व्याख्या आणि आपल्या !!{derivation} पद्धतींचे नियम देता यावेत,
यासाठी प्रथम-क्रम तर्कशास्त्रातील !!a{sentence} म्हणून नेमके काय गणले
जाते याची काटेकोर व्याख्या सर्वप्रथम द्यावी लागेल. मूलभूत कल्पना समजणे
पुरेसे सोपे आहे: फक्त !!{predicate}s आणि !!{constant}s यांपासून आपण
$\Atom{\Obj P}{\Obj a}$ यासारखी काही साधी !!{sentence}s रचू शकतो. मग
संयोजक आणि संख्यापक वापरून त्यांपासून अधिक जटिल वाक्ये रचतो. पण अधिक
जटिल !!{sentence}s रचण्याची अनुमती देणारे नियम नेमके कोणते? ते स्पष्ट
करावेच लागतील; अन्यथा ``प्रथम-क्रम तर्कशास्त्रातील !!{sentence}'' याची
आपण पुरेशी काटेकोर व्याख्या केलेली नाही. येथे काही अडचणी आहेत. पहिली
अडचण म्हणजे योग्य चिन्हमालाच !!{sentence}s म्हणून गणल्या जातील याची
खात्री करणे. दुसरी म्हणजे हे अशा रीतीने करणे की
\emph{सर्व}~!!{sentence}s विषयी गणितीय सिद्धता देता येतील. शेवटी,
``$!A$ मधील प्रत्येक $\Obj x$ च्या जागी~$\Obj b$ ठेवा'' यासारख्या
!!{sentence}s वरील काही प्राथमिक क्रियांच्याही काटेकोर व्याख्या द्याव्या
लागतील. प्रथम-क्रम तर्कशास्त्रातील संख्यापक आणि !!{variable}s यांमुळे हे
कसे करावे हे पूर्णपणे स्पष्ट नसते, ही अडचण आहे. उदाहरणार्थ,
$\lexists[\Obj x][\Atom{\Obj P}{\Obj a}]$ हे !!a{sentence} म्हणून
गणले जावे का? $\lexists[\Obj x][\lexists[\Obj x][\Atom{\Obj P}{\Obj
x}]]$ चे काय? आणि ``$(\Atom{\Obj P}{\Obj x} \land \lexists[\Obj
x][\Atom{\Obj P}{\Obj x}])$ मध्ये $\Obj x$ च्या जागी~$\Obj b$ ठेवा''
याचा परिणाम काय असावा?

\end{document}
